% GENERATED FILE. Edit canonical lessons, then run npm run generate:books.
% canonical-source-hash: fnv1a64:d81c77ad9b64538d
% canonical-lessons: JA-C35-satou, JA-C35-udon, JA-C35-sushi, JA-C35-asagohan, JA-C35-okashi

\chapter{Sugar, Udon noodles, Sushi, Breakfast, Sweets}
\label{ch:ja-sugar-udon-noodles-sushi-breakfast-sweets}
\hlchaptermodality{35}

\begin{chapteropening}
\textbf{By the end of this chapter:} \emph{I can say sugar, udon noodles, sushi, breakfast and sweets, snacks.}

Complete the last lesson of chapter 35: i can say sugar, udon noodles, sushi, breakfast and sweets, snacks.
\end{chapteropening}

\section[satou]{\ja{さとう} (satou) — sugar}

\label{lesson:JA-C35-satou}



\begin{quote}
Before the new one: say the Japanese for tofu, then the Japanese for salt.
\end{quote}

\subsection*{You'll want to know: \ja{さとう}}
\textbf{\ja{さとう}} — \emph{satou} — \textquotedblleft{}sugar\textquotedblright{}.

\begin{grammarlens}[title={What you've built}]
One more word for this chapter.
\end{grammarlens}

\subsection*{Your turn}
\begin{itemize}
\raggedright
  \item \emph{Say it:} \emph{satou}
  \item \emph{Say it:} \emph{satou}, once more
  \item \emph{Say it:} say \emph{satou}
  \item \emph{Recall:} say the Japanese for tofu, then the Japanese for salt, then say \emph{satou} again
\end{itemize}

\begin{tcolorbox}[breakable,colback=teal!4,colframe=teal!35!black,title={Before you move on}]
What does \ja{さとう} mean? (\textquotedblleft{}Sugar\textquotedblright{}.) Say it once more.
\end{tcolorbox}
\section[udon]{\ja{うどん} (udon) — udon noodles}

\label{lesson:JA-C35-udon}



\begin{quote}
Before the new one: say the Japanese for salt, then the Japanese for sugar.
\end{quote}

\subsection*{You'll want to know: \ja{うどん}}
\textbf{\ja{うどん}} — \emph{udon} — \textquotedblleft{}udon noodles\textquotedblright{}.

\begin{grammarlens}[title={What you've built}]
One more word for this chapter.
\end{grammarlens}

\subsection*{Your turn}
\begin{itemize}
\raggedright
  \item \emph{Say it:} \emph{udon}
  \item \emph{Say it:} \emph{udon}, once more
  \item \emph{Say it:} say \emph{udon}
  \item \emph{Recall:} say the Japanese for salt, then the Japanese for sugar, then say \emph{udon} again
\end{itemize}

\begin{tcolorbox}[breakable,colback=teal!4,colframe=teal!35!black,title={Before you move on}]
What does \ja{うどん} mean? (\textquotedblleft{}Udon noodles\textquotedblright{}.) Say it once more.
\end{tcolorbox}
\section[sushi]{\ja{すし} (sushi) — sushi}

\label{lesson:JA-C35-sushi}



\begin{quote}
Before the new one: say the Japanese for sugar, then the Japanese for udon noodles.
\end{quote}

\subsection*{You'll want to know: \ja{すし}}
\textbf{\ja{すし}} — \emph{sushi} — \textquotedblleft{}sushi\textquotedblright{}.

\begin{grammarlens}[title={What you've built}]
One more word for this chapter.
\end{grammarlens}

\subsection*{Your turn}
\begin{itemize}
\raggedright
  \item \emph{Say it:} \emph{sushi}
  \item \emph{Say it:} \emph{sushi}, once more
  \item \emph{Say it:} say \emph{sushi}
  \item \emph{Recall:} say the Japanese for sugar, then the Japanese for udon noodles, then say \emph{sushi} again
\end{itemize}

\begin{tcolorbox}[breakable,colback=teal!4,colframe=teal!35!black,title={Before you move on}]
What does \ja{すし} mean? (\textquotedblleft{}Sushi\textquotedblright{}.) Say it once more.
\end{tcolorbox}
\section[asagohan]{\ja{あさごはん} (asagohan) — breakfast}

\label{lesson:JA-C35-asagohan}



\begin{quote}
Before the new one: say the Japanese for udon noodles, then the Japanese for sushi.
\end{quote}

\subsection*{You'll want to know: \ja{あさごはん}}
\textbf{\ja{あさごはん}} — \emph{asagohan} — \textquotedblleft{}breakfast\textquotedblright{}.

\begin{grammarlens}[title={What you've built}]
One more word for this chapter.
\end{grammarlens}

\subsection*{Your turn}
\begin{itemize}
\raggedright
  \item \emph{Say it:} \emph{asagohan}
  \item \emph{Say it:} \emph{asagohan}, once more
  \item \emph{Say it:} say \emph{asagohan}
  \item \emph{Recall:} say the Japanese for udon noodles, then the Japanese for sushi, then say \emph{asagohan} again
\end{itemize}

\begin{tcolorbox}[breakable,colback=teal!4,colframe=teal!35!black,title={Before you move on}]
What does \ja{あさごはん} mean? (\textquotedblleft{}Breakfast\textquotedblright{}.) Say it once more.
\end{tcolorbox}
\section[okashi]{\ja{おかし} (okashi) — sweets, snacks}

\label{lesson:JA-C35-okashi}



\begin{quote}
Before the new one: say the Japanese for sushi, then the Japanese for breakfast.
\end{quote}

\subsection*{You'll want to know: \ja{おかし}}
\textbf{\ja{おかし}} — \emph{okashi} — \textquotedblleft{}sweets, snacks\textquotedblright{}.

\begin{grammarlens}[title={What you've built}]
One more word for this chapter.
\end{grammarlens}

\subsection*{Your turn}
\begin{itemize}
\raggedright
  \item \emph{Say it:} \emph{okashi}
  \item \emph{Say it:} \emph{okashi}, once more
  \item \emph{Say it:} say \emph{okashi}
  \item \emph{Recall:} say the Japanese for sushi, then the Japanese for breakfast, then say \emph{okashi} again
\end{itemize}

\begin{tcolorbox}[breakable,colback=teal!4,colframe=teal!35!black,title={Before you move on}]
What does \ja{おかし} mean? (\textquotedblleft{}Sweets, snacks\textquotedblright{}.) Say it once more.
\end{tcolorbox}
